\section{总结与延伸}

这节课最重要的贡献不在于公开某个神秘配方，而在于把 Agentic 模型训练的真实重心讲清楚了：\textbf{数据与 grader 是能力上限，系统效率是迭代速度，二者共同决定最终模型质量。}

\begin{longtable}{p{2.1cm}p{3.6cm}p{5.2cm}}
\toprule
\textbf{主题} & \textbf{课程关键观点} & \textbf{落地建议} \\
\midrule
训练目标 & 端到端质量提升依赖训练闭环，而非单点 prompt & 以任务轨迹为单位定义优化目标 \\
数据策略 & 高质量合成与定向补短板比盲目扩量更有效 & 建立任务分桶和难度课程机制 \\
Grader 体系 & 评分器定义得当，等于解决一半问题 & 构建 verifier + rubric + judge 组合 \\
反作弊 & 模型会主动 exploit grader 漏洞 & 将 anti-cheat 设计前置到环境与评估 \\
效率工程 & 异步采样与系统吞吐决定研发节奏 & 先优化迭代效率，再追最优配置 \\
工具泛化 & MCP 场景下需适配未知工具 & 强化 instruction following 与反思纠错 \\
LoRA in RL & 价值不仅是省算力，更是正则化稳定性 & 探索期优先 LoRA，冲刺期再放宽 \\
未来方向 & test-time scaling + self-improvement & 建立统一评估预算与持续回流机制 \\
\bottomrule
\end{longtable}

\begin{importantbox}{一句话收束}
\textit{``训练 Agent 不只是把模型训得更聪明，而是把模型训得更可控、更可评估、并且在真实约束下持续可进化。''}
\end{importantbox}

\subsection*{延伸阅读}
\begin{enumerate}
    \item OpenAI, Anthropic, DeepMind 近两年关于 process supervision 与 scalable oversight 的论文，用于补充 ``过程可评估'' 视角。
    \item SWE-bench 系列与相关 agent benchmark 报告，重点关注 ``quality vs cost'' 联合指标而非单一 pass rate。
    \item LoRA without Regret 等参数高效微调研究，理解其在 RL 稳定性上的潜在机制。
    \item 关于 asynchronous RL 与 distributed training 的系统论文，补足 sampler/trainer 解耦设计细节。
    \item 工业实践博客与技术分享（MCP/tool-use、sandbox security、eval platform），用于建立端到端训练工程认知。
\end{enumerate}

\end{document}
